\chapter*{宣言事項}
\markboth{宣言事項}{宣言事項}
\addcontentsline{toc}{chapter}{宣言事項}

\section*{AIツールの使用}

本プロジェクトにおいて、Claude（Anthropic、\url{https://claude.ai/}）およびOpenAIのGPTモデル（\texttt{GPT-4.1-mini}、\texttt{GPT-4o-mini}、\texttt{GPT-4o}、\texttt{GPT-4.5}を含む）（OpenAI、\url{https://platform.openai.com/}）を使用したことを申告する。Claudeは本報告書の一部へのフィードバック提供、Pythonコードのデバッグ支援、およびプロジェクト遂行中に読んだ論文の理解深化のために使用した。本報告書中のいくつかの図は、最初に筆者が手描きし、その後ChatGPT（\texttt{GPT-4.5}）のウェブインターフェースを通じてクリーンなベクタースタイルの図へと精緻化した。\par

GPTモデルは、第\ref{chap:development1_emo_steering}章で説明する感情書き換えデータセットの生成に使用し、プルチックの8つの基本感情にわたって3段階の強度で発話の感情的バリアントを生成した。また、ステアリング実験の出力評価にも使用し、抽出されたベクトルに沿った介入によって生じた感情強度とテクスチャの変化を評価した。この使用は方法論（第\ref{sec:steering_experiments}節、第\ref{sec:ch2-causal-validation}節）に明示的に記載されている。AIが生成したコンテンツを自身の成果として提示したことはない。\par

\section*{倫理的考慮}

本プロジェクトは従来の意味での人間被験者を含まないが、参加者が生成テキストの感情的特質を評価した人間評価研究を含んでいる。参加は任意であり、個人を特定できる情報は収集されず、すべての参加者のインフォームドコンセントを得た上で研究を実施した。\par

より広い観点から、本プロジェクトは言語モデルにおける感情操作を解釈可能かつ制御可能にするという目標に貢献するものである。生成テキストの感情的トーンを制御するシステムは、例えば大規模な操作的または欺瞞的コンテンツの生成といった潜在的な悪用に関する問いを提起する。本研究で開発した手法は解釈可能な幾何学的成分のレベルで動作しており、この特性は難読化よりも監査可能性と透明性を支持するものである。それでもなお、感情ステアリング技術の二重使用の可能性は、そのようなシステムが成熟するにつれて継続的な注意を要する。\par

\section*{持続可能性}

本プロジェクトの計算コストは小規模なものであった。訓練は一切行っておらず、すなわちすべての実験において凍結された事前学習済みモデルの重みを使用した。活性化抽出とステアリング実験は単一のマシン上で実行し、小規模モデルにはCPU推論を、\texttt{Llama-3.1-8B-Instruct}には民生用GPUを使用した。データセット構築ステップでは\texttt{GPT-4.1-mini}をAPI経由で使用して多数の短い書き換えを生成した。モデルの訓練は行っていないが、このリモート計算がプロジェクトの主な外部計算コストとなっている。
\section*{データと資料}

活性化抽出、CAAベクトル構築、分解、プールドPCA、ステアリング実験、および人間評価分析のすべてのコードは、本プロジェクトの正式リポジトリで管理されている：

\begin{itemize}
    \item \url{https://github.com/maplesugano/EmotionEngine_v2} --- 本報告書の正式リポジトリ。すべての最終実験、分析スクリプト、および再現性チェックリスト（付録~\ref{app:reproducibility}）を含む
\end{itemize}

以下の2つの旧リポジトリは歴史的参照のためにのみ存在し、本報告書で報告されたいかなる結果の再現にも必要ない：
\begin{itemize}
    \item \url{https://github.com/maplesugano/EmotionEngine} --- 現在のアーキテクチャに先行する初期段階のプロトタイプ
    \item \url{https://github.com/maplesugano/Dissertation} --- プロジェクト開始初週の初期メモと探索的コード
\end{itemize}
